% Extremalprobleme · Maximum bestimmen · Prüfungs-Fokus Abitur GK – bank/extremalprobleme/e3.jsonl (zusammenbau v0.6, Prüfungs-Fokus)
\einheitenkopf[e3]{Einheit 3 · Maximum bestimmen und deuten}
\zweigzeile{Abi GK ×7}
\setcounter{aufgabe}{0}

% Anlauf (ohne Prüfkennung)
\begin{aufgabe}{Maximum bestimmen – Anlauf}
\begin{teile}
\teil Frage: „Wie groß ist der größtmögliche Flächeninhalt?“ Was ist verlangt? \kreuz{die Stelle $x$} \\ \kreuz{die Seitenlängen} \\ \kreuz{der Flächeninhalt}
\teil Die Zielfunktion $A(a) = 30a - 2a^2$, $0 < a < 15$, beschreibt einen Flächeninhalt. Bestimme die Stelle mit $A'(a) = 0$ im Definitionsbereich. a = \leerfeld
\teil Für ein Beet an einer Mauer mit 24 m Zaun für drei Seiten gilt $A(a) = 24a - 2a^2$ ($a$ senkrecht zur Mauer, in m). Bestimme Maße und Flächeninhalt des größten Beets und weise nach, dass es ein Maximum ist.
\end{teile}
\end{aufgabe}

\begin{aufgabe}{Maximum bestimmen – Prüfungsaufgaben}
\begin{teile}
\teil Der Ursprung und $P(u | f(u))$ mit $u > 0$ auf dem Graphen von $f(x) = (x + 3) \cdot e^{-0{,}5x}$ sind gegenüberliegende Ecken eines achsenparallelen Rechtecks; genau ein $u$ liefert den größten Flächeninhalt. Bestimme die Koordinaten von $P$ für diesen Fall (Kontrolle: $A'(u) = (-0{,}5u^2 + 0{,}5u + 3) \cdot e^{-0{,}5u}$). \hfill (A22)
\teil Der Ursprung und $P(u | f(u))$ mit $u > 0$ auf dem Graphen von $f(x) = (x + 4) \cdot e^{-x}$ sind gegenüberliegende Ecken eines achsenparallelen Rechtecks; genau ein $u$ liefert den größten Flächeninhalt. Bestimme die Koordinaten von $P$ für diesen Fall (Kontrolle: $A'(u) = (-u^2 - 2u + 4) \cdot e^{-u}$). \hfill (A22)
\teil Der Ursprung und $P(u | f(u))$ mit $0 < u < 3$ auf dem Graphen von $f(x) = 9 - x^2$ sind gegenüberliegende Ecken eines achsenparallelen Rechtecks. Bestimme die Koordinaten von $P$, für die der Flächeninhalt am größten ist. \hfill (A22)
\teil Ein Skispringer fliegt entlang $f(x) = -0{,}01x^2 + 40$, der Hang verläuft entlang $g(x) = -0{,}5x + 40$ (1 LE = 1 m); der Sprung endet bei $x = 50$. Weise nach, dass der größte vertikale Abstand zwischen Springer und Hang $6{,}25$ m beträgt. Auf die hinreichende Bedingung darf verzichtet werden. \hfill (A18)
\teil Die Dicke eines Bauteils wird auf $[0; 4{,}5]$ durch $d(x) = f(x) - h(x)$ mit $f(x) = -0{,}1x^3 + 0{,}3x^2 + 0{,}9x + 1$ und $h(x) = 1$ beschrieben ($x$ und $d$ in dm). Sie darf $2{,}8$ dm nicht überschreiten. Prüfe, ob die Vorgabe eingehalten ist. \hfill (A21)
\teil Die Funktion $f(x) = x^4 - 5x^2 + 4$ wird für $-1 \le x \le 1$ durch $p(x) = -4x^2 + 4$ angenähert; $p$ liegt dort nicht unter $f$. Weise nach, dass die größte Abweichung $\tfrac{1}{4}$ beträgt. \hfill (A24)
\teil Unter dem Graphen von $f(x) = -x^2 + 10x$ liegen die Trapeze mit den Ecken $(0 | 0)$, $(10 | 0)$, $(10 - u | f(u))$ und $(u | f(u))$, $0 < u < 5$; ihr Flächeninhalt ist $T(u) = (10 - u) \cdot f(u)$. Eines hat den größten Inhalt. Bestimme den zugehörigen Wert von $u$. \hfill (A19) \\ u = \leerfeld
\teil Unter dem Graphen von $f(x) = -0{,}5x^2 + 3x$ liegen die Trapeze mit den Ecken $(0 | 0)$, $(6 | 0)$, $(6 - u | f(u))$ und $(u | f(u))$, $0 < u < 3$; ihr Flächeninhalt ist $T(u) = (6 - u) \cdot f(u)$. Eines hat den größten Inhalt. Bestimme den zugehörigen Wert von $u$ und den Inhalt. \hfill (A19) \\ u = \leerfeld
\teil Unter dem Graphen von $f(x) = -2x^2 + 8x$ liegen die Trapeze mit den Ecken $(0 | 0)$, $(4 | 0)$, $(4 - u | f(u))$ und $(u | f(u))$, $0 < u < 2$; ihr Flächeninhalt ist $T(u) = (4 - u) \cdot f(u)$. Eines hat den größten Inhalt. Bestimme den zugehörigen Wert von $u$. \hfill (A19) \\ u = \leerfeld
\teil Für $0 < x < 2$ legt der Punkt $P(x | f(x))$ auf dem Graphen von $f(x) = 4 - x^2$ mit dem Ursprung ein achsenparalleles Rechteck fest, dessen Flächeninhalt für genau ein $x$ maximal ist. Weise nach, dass diese Stelle nicht $x = 1$ ist. \hfill (A21)
\teil Für $0 < x < 3$ legt der Punkt $P(x | f(x))$ auf dem Graphen von $f(x) = 6 - 2x$ mit dem Ursprung ein achsenparalleles Rechteck fest, dessen Flächeninhalt für genau ein $x$ maximal ist. Weise nach, dass diese Stelle nicht $x = 2$ ist. \hfill (A21)
\teil Für $0 < x < 2$ legt der Punkt $P(x | f(x))$ auf dem Graphen von $f(x) = x^3 - 5x^2 + 6x + 1$ mit dem Ursprung ein achsenparalleles Rechteck fest, dessen Flächeninhalt für genau ein $x$ maximal ist. Weise nach, dass diese Stelle nicht $x = 1$ ist. \hfill (A21)
\end{teile}
\end{aufgabe}

\medskip
\uebersichtskasten{Rechnen: Zielfunktion ableiten, A'(x) = 0 lösen, Lösungen außerhalb des Definitionsbereichs verwerfen; Art mit A''(x) oder dem Vorzeichenwechsel von A'; auf abgeschlossenen Bereichen die Randwerte vergleichen. \\ A(a) = a · (ℓ − 2a) = ℓa − 2a²: A'(a) = ℓ − 4a = 0 ⇔ a = ℓ/4; A''(a) = −4 < 0 → Maximum; b = ℓ − 2 · ℓ/4 = ℓ/2. \\ Antwort vollständig: angeben, was gefragt ist – Stelle, Seitenlängen und Extremwert mit Einheit; die Extremstelle allein ist kein Flächeninhalt. \\ Existenz ohne Rechnung: An den Rändern des Bereichs entartet die Figur (Breite oder Höhe gegen null) – dazwischen liegt ein größtes Exemplar; ein kleinstes gibt es nicht, wenn die Randfälle ausgeschlossen sind. \\ Abstand zweier Graphen: der vertikale Abstand ist die Differenz der Funktionswerte – Zielfunktion d(x) = f(x) − g(x) (oberer minus unterer Graph), dann wie oben.}
